% ============================================================================
% WRAP-UP
% ============================================================================

\begin{frame}{Debrief}

\textbf{Three screens.} A completeness check that passed, or one that did not. A confusion matrix. One disagreement where the model was right and the coder wrong.

\vspace{0.3cm}

\textbf{From the shared document.} Your agreement with the coders against the models'. The range of $\kappa$ across the room: same data, same code, different runs.

\vspace{0.3cm}

\textbf{One question.} Would you put today's labels in a paper? What would the referee ask first?

\vspace{0.3cm}

\textbf{Before you leave.} Download your bundle (Part 7). In your project: prompts 10 and 11, the section ``APIs, keys and LLM annotation'' of \texttt{template/CLAUDE.md}, then prompt 05 and the commit.

\end{frame}

\begin{frame}{Next Session: 18 November 2026}

\textbf{Session 4: Extracting Information, With and Without an LLM}

\vspace{0.2cm}

\begin{itemize}
    \item The classic toolbox first: dictionaries, embeddings, a supervised classifier
    \item Then LLM extraction: from a text to a table (who, what, when, how much)
    \item Prompt 11 on your own corpus: bring 100 texts you have coded yourself
    \item Keep working in your folder. End every session with prompt 05.
\end{itemize}

\vspace{0.3cm}

\begin{alertblock}{Carry-over question}
The model agrees with the human coders on 75\% of the sentences. Your regression uses its labels. What would you have to show for a referee to accept the coefficient?
\end{alertblock}

\end{frame}

% ============================================================================
% REFERENCES
% ============================================================================

\begin{frame}{References (1/2)}

{\scriptsize

Barrie, C., Palaiologou, E.\ and T\"ornberg, P.\ (2026). Prompt Stability Scoring for Text Annotation with Large Language Models. arXiv:2407.02039.

\vspace{0.08cm}

Benoit, K., De Marchi, S., Laver, C., Laver, M.\ and Ma, J.\ (2026). Using Large Language Models to Analyze Political Texts through Natural Language Understanding. \textit{American Journal of Political Science} 70(4), 1589--1605.

\vspace{0.08cm}

Egami, N., Hinck, M., Stewart, B.\ M.\ and Wei, H.\ (2023). Using Imperfect Surrogates for Downstream Inference: Design-based Supervised Learning for Social Science Applications of Large Language Models. \textit{NeurIPS} 36.

\vspace{0.08cm}

Gilardi, F., Alizadeh, M.\ and Kubli, M.\ (2023). ChatGPT Outperforms Crowd Workers for Text-Annotation Tasks. \textit{PNAS} 120(30), e2305016120.

\vspace{0.08cm}

Halterman, A.\ and Keith, K.\ A.\ (2026). Codebook LLMs: Evaluating LLMs as Measurement Tools for Political Science Concepts. \textit{Political Analysis} 34(2), 188--204.

\vspace{0.08cm}

%%MANIFESTOREFS%%

\vspace{0.08cm}

Ornstein, J.\ T., Blasingame, E.\ N.\ and Truscott, J.\ S.\ (2025). How to Train Your Stochastic Parrot: Large Language Models for Political Texts. \textit{Political Science Research and Methods} 13(2), 264--281.

}

\end{frame}

\begin{frame}{References (2/2)}

{\scriptsize

Palmer, A., Smith, N.\ A.\ and Spirling, A.\ (2024). Using Proprietary Language Models in Academic Research Requires Explicit Justification. \textit{Nature Computational Science} 4(1), 2--3.

\vspace{0.08cm}

Reiss, M.\ V.\ (2023). Testing the Reliability of ChatGPT for Text Annotation and Classification: A Cautionary Remark. arXiv:2304.11085.

\vspace{0.08cm}

T\"ornberg, P.\ (2024). Best Practices for Text Annotation with Large Language Models. \textit{Sociologica} 18(2), 67--85.

\vspace{0.08cm}

T\"ornberg, P.\ (2025). Large Language Models Outperform Expert Coders and Supervised Classifiers at Annotating Political Social Media Messages. \textit{Social Science Computer Review} 43(6), 1181--1195.

\vspace{0.08cm}

Ziems, C., Held, W., Shaikh, O., Chen, J., Zhang, Z.\ and Yang, D.\ (2024). Can Large Language Models Transform Computational Social Science? \textit{Computational Linguistics} 50(1), 237--291.

\vspace{0.15cm}

\textbf{Documentation} (consulted 3 November 2026): Anthropic, Claude API (models, pricing, structured outputs, effort); OpenAI, Google AI, Mistral AI pricing pages; Ollama documentation; Science IT, ScienceCluster documentation (\href{https://docs.s3it.uzh.ch}{\textcolor{titleblue}{docs.s3it.uzh.ch}}); Claude Code setup and authentication (\href{https://code.claude.com/docs}{\textcolor{titleblue}{code.claude.com/docs}}).

}

\end{frame}
